\documentclass[12pt]{jlreq}
\usepackage{amsmath, amssymb}

\newcommand{\C}{\mathbb{C}}
\newcommand{\R}{\mathbb{R}}
\newcommand{\Z}{\mathbb{Z}}
\newcommand{\Tr}{\operatorname{Tr}}
\newcommand{\Impart}{\operatorname{Im}}
\newcommand{\AF}{\operatorname{AF}}
\newcommand{\EG}{\operatorname{EG}}
\newcommand{\AX}{\operatorname{AX}}
\newcommand{\GL}{\operatorname{GL}}
\newcommand{\Hom}{\operatorname{Hom}}
\newcommand{\Ker}{\operatorname{Ker}}
\newcommand{\Id}{\operatorname{Id}}
\newcommand{\Sym}{\operatorname{Sym}}

\begin{document}

$A = \mathbb{C}[T]$ を複素数体 $\mathbb{C}$ 上の $1$ 変数多項式環とし，$K = \mathbb{C}(T)$ をその分体とする。$B$ を $\mathbb{C}$ 上の $3$ 変数多項式環 $\mathbb{C}[X,Y,Z]$ のイデアル $I = (XY, YZ, ZX)$ による商環 $\mathbb{C}[X,Y,Z]/I$ とする。また $\mathbb{C}$ 上の環準同型 $f: A \to B$ を $f(T) = \overline{X+Y+Z}$ で定める。ただしここで，$\overline{X+Y+Z}$ は $B$ における $X+Y+Z$ の同値類である。以下の問に答えよ。

\begin{enumerate}
  \item $B$ を $f$ により $A$ 加群とみるとき，$B$ は自由 $A$ 加群であることを示し，その階数を求めよ。
  \item テンソル積 $B \otimes_A K$ のべき等元をすべて求めよ。(ただしここで，$x \in B \otimes_A K$ がべき等元であるとは，$x^2 = x$ を満たすことである。)
  \item 包含写像 $A \to K$ がひきおこす単射 $B \to B \otimes_A K$ により $B$ を $B \otimes_A K$ の部分環と考え，$B \otimes_A K$ のすべてのべき等元で $B$ 上生成される $B \otimes_A K$ の部分環を $C$ で表す。$C/B$ の $\mathbb{C}$ 上の線形空間としての次元を求めよ。
  \item 乗法群の包含写像 $B^\times \to C^\times$ が同型かどうか判定せよ。
\end{enumerate}

\end{document}